% ------------------------------------------------------------
% LaTeX Template für die DHBW zum Schnellstart!
% Original: https://github.wdf.sap.corp/vtgermany/LaTeX-Template-DHBW
% ------------------------------------------------------------
% ---- Präambel mit Angaben zum Dokument
\documentclass[
	fontsize=11pt,           % Leitlinien sprechen von Schriftgröße 12.
	paper=A4,
	twoside=false,
	listof=totoc,            % Tabellen- und Abbildungsverzeichnis ins Inhaltsverzeichnis
	bibliography=totoc,      % Literaturverzeichnis ins Inhaltsverzeichnis aufnehmen
	titlepage,               % Titlepage-Umgebung anstatt \maketitle
	headsepline,             % horizontale Linie unter Kolumnentitel
	abstract,              % Überschrift einschalten, Abstract muss in {abstract}-Umgebung stehen
]{scrreprt}                  % Verwendung von KOMA-Report
\usepackage[utf8]{inputenc}  % UTF8 Encoding einschalten
\usepackage[ngerman]{babel}  % Neue deutsche Rechtschreibung
\usepackage[T1]{fontenc}     % Ausgabe von westeuropäischen Zeichen (auch Umlaute)
\usepackage{microtype}       % Trennung von Wörtern wird besser umgesetzt
\usepackage{lmodern}         % Nicht-gerasterte Schriftarten (bei MikTeX erforderlich)
\usepackage{graphicx}        % Einbinden von Grafiken erlauben
\usepackage{subcaption}
\usepackage{wrapfig}         % Grafiken fließend im Text
\usepackage{setspace}        % Zeilenabstand \singlespacing, \onehalfspaceing, \doublespacing
\usepackage[
	%showframe,                % Ränder anzeigen lassen
	left=2cm, right=2cm,
	top=2.5cm,  bottom=2.5cm,
	includeheadfoot
]{geometry}                      % Seitenlayout einstellen
\usepackage{scrlayer-scrpage}    % Gestaltung von Fuß- und Kopfzeilen
\usepackage{acronym}             % Abkürzungen, Abkürzungsverzeichnis
\usepackage{titletoc}            % Anpassungen am Inhaltsverzeichnis
\usepackage{longtable}
\contentsmargin{0.75cm}          % Abstand im Inhaltsverzeichnis zw. Punkt und Seitenzahl
\usepackage[                     % Klickbare Links (enth. auch "nameref", "url" Package)
  hidelinks,                     % Blende die "URL Boxen" aus.
  breaklinks=true                % Breche zu lange URLs am Zeilenende um
]{hyperref}
\usepackage[hypcap=true, font=small]{caption}% Anker Anpassung für Referenzen
\urlstyle{same}                  % Aktuelle Schrift auch für URLs
% Anpassung von autoref für Gleichungen (ergänzt runde Klammern) und Algorithm.
% Anstatt "Listing" kann auch z.B. "Code-Ausschnitt" verwendet werden. Dies sollte
% jedoch synchron gehalten werden mit \lstlistingname (siehe weiter unten).
\addto\extrasngerman{%
	\def\equationautorefname~#1\null{Gleichung~(#1)\null}
	\def\lstnumberautorefname{Zeile}
	\def\lstlistingautorefname{Listing}
	\def\algorithmautorefname{Algorithmus}
	% Damit einheitlich "Abschnitt 1.2[.3]" verwendet wird und nicht "Unterabschnitt 1.2.3"
	% \def\subsectionautorefname{Abschnitt}
}

% ---- Abstand verkleinern von der Überschrift 
\renewcommand*{\chapterheadstartvskip}{\vspace*{.5\baselineskip}}

% Hierdurch werden Schusterjungen und Hurenkinder vermieden, d.h. einzelne Wörter
% auf der nächsten Seite oder in einer einzigen Zeile.
% LaTeX kann diese dennoch erzeugen, falls das Layout ansonsten nicht umsetzbar ist.
% Diese Werte sind aber gute Startwerte.
\widowpenalty10000
\clubpenalty10000

% ---- Für das Quellenverzeichnis
\usepackage[
	backend = biber,                % Verweis auf biber
	language = auto,
	style = numeric,                % Nummerierung der Quellen mit Zahlen
	sorting = none,                 % none = Sortierung nach der Erscheinung im Dokument
	sortcites = true,               % Sortiert die Quellen innerhalb eines cite-Befehls
	block = space,                  % Extra Leerzeichen zwischen Blocks
	hyperref = true,                % Links sind klickbar auch in der Quelle
	%backref = true,                % Referenz, auf den Text an die zitierte Stelle
	bibencoding = auto,
	giveninits = true,              % Vornamen werden abgekürzt
	doi=false,                      % DOI nicht anzeigen
	isbn=false,                     % ISBN nicht anzeigen
    alldates=short                  % Datum immer als DD.MM.YYYY anzeigen
]{biblatex}
\addbibresource{Inhalt/literatur.bib}
\setcounter{biburlnumpenalty}{3000}     % Umbruchgrenze für Zahlen
\setcounter{biburlucpenalty}{6000}      % Umbruchgrenze für Großbuchstaben
\setcounter{biburllcpenalty}{9000}      % Umbruchgrenze für Kleinbuchstaben
\DeclareNameAlias{default}{family-given}  % Nachname vor dem Vornamen
\AtBeginBibliography{\renewcommand{\multinamedelim}{\addslash\space
}\renewcommand{\finalnamedelim}{\multinamedelim}}  % Schrägstrich zwischen den Autorennamen
\DefineBibliographyStrings{german}{
  urlseen = {Einsichtnahme:},                      % Ändern des Titels von "besucht am"
}
\usepackage[babel,german=quotes]{csquotes}         % Deutsche Anführungszeichen + Zitate


% ---- Für Mathevorlage
\usepackage{amsmath}    % Erweiterung vom Mathe-Satz
\usepackage{amssymb}    % Lädt amsfonts und weitere Symbole
\usepackage{MnSymbol}   % Für Symbole, die in amssymb nicht enthalten sind.


% ---- Für Quellcodevorlage
%\usepackage{scrhack}                    % Hack zur Verw. von listings in KOMA-Script
\usepackage{minted}
\setminted{linenos=true}
\newcommand\mynumberformat{\def\FancyVerbFormatLine##1{{\theFancyVerbLine} ##1}}
\RecustomVerbatimEnvironment{Verbatim}{BVerbatim}
{formatcom=\mynumberformat}
\usepackage{listings}                   % Darstellung von Quellcode
\usepackage[dvipsnames]{xcolor}
% -- Eigene Farben für den Quellcode
\definecolor{JavaLila}{rgb}{0.4,0.1,0.4}
\definecolor{JavaGruen}{rgb}{0.3,0.5,0.4}
\definecolor{JavaBlau}{rgb}{0.0,0.0,1.0}
\definecolor{ABAPKeywordsBlue}{HTML}{6000ff}
\definecolor{ABAPCommentGrey}{HTML}{808080}
\definecolor{ABAPStringGreen}{HTML}{4da619}
\definecolor{PyKeywordsBlue}{HTML}{0000AC}
\definecolor{PyCommentGrey}{HTML}{808080}
\definecolor{PyStringGreen}{HTML}{008080}
% -- Farben für ABAP CDS
\definecolor{CDSString}{HTML}{FF8C00}
\definecolor{CDSKeywords}{HTML}{6000ff}
\definecolor{CDSAnnotation}{HTML}{00BFFF}
\definecolor{CDSComment}{HTML}{808080}
\definecolor{CDSFunc}{HTML}{FF0000}

% -- Default Listing-Styles

\lstset{
	% Das Paket "listings" kann kein UTF-8. Deswegen werden hier 
	% die häufigsten Zeichen definiert (ä,ö,ü,...)
	literate=%
		{á}{{\'a}}1 {é}{{\'e}}1 {í}{{\'i}}1 {ó}{{\'o}}1 {ú}{{\'u}}1
		{Á}{{\'A}}1 {É}{{\'E}}1 {Í}{{\'I}}1 {Ó}{{\'O}}1 {Ú}{{\'U}}1
		{à}{{\`a}}1 {è}{{\`e}}1 {ì}{{\`i}}1 {ò}{{\`o}}1 {ù}{{\`u}}1
		{À}{{\`A}}1 {È}{{\'E}}1 {Ì}{{\`I}}1 {Ò}{{\`O}}1 {Ù}{{\`U}}1
		{ä}{{\"a}}1 {ë}{{\"e}}1 {ï}{{\"i}}1 {ö}{{\"o}}1 {ü}{{\"u}}1
		{Ä}{{\"A}}1 {Ë}{{\"E}}1 {Ï}{{\"I}}1 {Ö}{{\"O}}1 {Ü}{{\"U}}1
		{â}{{\^a}}1 {ê}{{\^e}}1 {î}{{\^i}}1 {ô}{{\^o}}1 {û}{{\^u}}1
		{Â}{{\^A}}1 {Ê}{{\^E}}1 {Î}{{\^I}}1 {Ô}{{\^O}}1 {Û}{{\^U}}1
		{œ}{{\oe}}1 {Œ}{{\OE}}1 {æ}{{\ae}}1 {Æ}{{\AE}}1 {ß}{{\ss}}1
		{ű}{{\H{u}}}1 {Ű}{{\H{U}}}1 {ő}{{\H{o}}}1 {Ő}{{\H{O}}}1
		{ç}{{\c c}}1 {Ç}{{\c C}}1 {ø}{{\o}}1 {å}{{\r a}}1 {Å}{{\r A}}1
		{€}{{\euro}}1 {£}{{\pounds}}1 {«}{{\guillemotleft}}1
		{»}{{\guillemotright}}1 {ñ}{{\~n}}1 {Ñ}{{\~N}}1 {¿}{{?`}}1,
	breaklines=true,        % Breche lange Zeilen um 
	breakatwhitespace=true, % Wenn möglich, bei Leerzeichen umbrechen
	% Symbol für Zeilenumbruch einfügen
	prebreak=\raisebox{0ex}[0ex][0ex]{\ensuremath{\rhookswarrow}},
	postbreak=\raisebox{0ex}[0ex][0ex]{\ensuremath{\rcurvearrowse\space}},
	tabsize=4,                                 % Setze die Breite eines Tabs
	basicstyle=\ttfamily\small,                % Grundsätzlicher Schriftstyle
	columns=fixed,                             % Besseres Schriftbild
	numbers=left,                              % Nummerierung der Zeilen
	%frame=single,                             % Umrandung des Codes
	showstringspaces=false,                    % Keine Leerzeichen hervorheben
	keywordstyle=\color{blue},
	ndkeywordstyle=\bfseries\color{darkgray},
	identifierstyle=\color{black},
	commentstyle=\itshape\color{JavaGruen},   % Kommentare in eigener Farbe
	stringstyle=\color{JavaBlau},             % Strings in eigener Farbe,
	captionpos=b,                             % Bild*unter*schrift
	xleftmargin=5.0ex
}

% ---- Eigener JAVA-Style für den Quellcode
\renewcommand{\ttdefault}{pcr}               % Schriftart, welche auch fett beinhaltet
\lstdefinestyle{EigenerJavaStyle}{
	language=Java,                             % Syntax Highlighting für Java
	%frame=single,                             % Umrandung des Codes
	keywordstyle=\bfseries\color{JavaLila},    % Keywords in eigener Farbe und fett
	commentstyle=\itshape\color{JavaGruen},    % Kommentare in eigener Farbe und italic
	stringstyle=\color{JavaBlau}               % Strings in eigener Farbe
}

% ---- Eigener ABAP-Style für den Quellcode
\renewcommand{\ttdefault}{pcr}
\lstdefinestyle{EigenerABAPStyle}{
	language=[R/3 6.10]ABAP,
	morestring=[b]\|,                          % Für Pipe-Strings
	morestring=[b]\`,                          % für Backtick-Strings
	keywordstyle=\bfseries\color{ABAPKeywordsBlue},
	commentstyle=\itshape\color{ABAPCommentGrey},
	stringstyle=\color{ABAPStringGreen},
	tabsize=2,
	morekeywords={
		types,
		@data,
		as,
		lower,
		start,
		selection,
		order,
		by,
		inner,
		join,
		key,
		end,
		cast
	}
}

% ---- Eigener Python-Style für den Quellcode
\renewcommand{\ttdefault}{pcr}
\lstdefinestyle{EigenerPythonStyle}{
	language=Python,
	columns=flexible,
	keywordstyle=\bfseries\color{PyKeywordsBlue},
	commentstyle=\itshape\color{PyCommentGrey},
	stringstyle=\color{PyStringGreen}
}

%----- ABAP-CDS-View language
\lstdefinelanguage{ABAPCDS}{
	sensitive=false,
	%Keywords
	morekeywords={define,
		view,
		as,
		select,
		from,
		inner,
		join,
		on,
		key,
		case,
		when,
		then,
		else,
		end,
		true,
		false,
		cast,
		where,
		and,
		distinct,
		group,
		by,
		having,
		min,
		sum,
		max,
		count,
		avg
	},
	%Methoden
	morekeywords=[2]{
		div,
		currency\_conversion,
		dats\_days\_between,
		concat\_with\_space,
		dats\_add_days,
		dats\_is\_valid,
		dats\_add\_months,
		unit\_conversion,
		division,
		mod,
		abs,
		floor,
		ceil,
		round,
		concat,
		replace,
		substring,
		left,
		right,
		length
	},
	morecomment=[s][\color{CDSAnnotation}]{@}{:},
	morecomment=[l][\itshape\color{CDSComment}]{//},
	morecomment=[s][\itshape\color{CDSComment}]{/*}{*/},
	morestring=[b][\color{CDSString}]',
	keywordstyle=\bfseries\color{CDSKeywords},
	keywordstyle=[2]\color{CDSFunc}
}

  % Weitere Details sind ausgelagert

\usepackage{algorithm}                  % Für Algorithmen-Umgebung (ähnlich wie lstlistings Umgebung)
\usepackage{algpseudocode}              % Für Pseudocode. Füge "[noend]" hinzu, wenn du kein "endif",
                                        % etc. haben willst.

\makeatletter                           % Sorgt dafür, dass man @ in Namen verwenden kann.
                                        % Ansonsten gibt es in der nächsten Zeile einen Compilefehler.
\renewcommand{\ALG@name}{Algorithmus}   % Umbenennen von "Algorithm" im Header der Listings.
\makeatother                            % Zeichen wieder zurücksetzen
\renewcommand{\lstlistingname}{Listing} % Erlaubt das Umbenennen von "Listing" in anderen Titel.

% ---- Tabellen
\usepackage{booktabs}  % Für schönere Tabellen. Enthält neue Befehle wie \midrule
\usepackage{multirow}  % Mehrzeilige Tabellen
\usepackage{siunitx}   % Für SI Einheiten und das Ausrichten Nachkommastellen
\sisetup{locale=DE, range-phrase={~bis~}, output-decimal-marker={,}} % Damit ein Komma und kein Punkt verwendet wird.
\usepackage{xfrac} % Für siunitx Option "fraction-function=\sfrac"

% ---- Für Definitionsboxen in der Einleitung
\usepackage{amsthm}                     % Liefert die Grundlagen für Theoreme
\usepackage[framemethod=tikz]{mdframed} % Boxen für die Umrandung
% ---- Definition für Highlight Boxen

% ---- Grundsätzliche Definition zum Style
\newtheoremstyle{defi}
  {\topsep}         % Abstand oben
  {\topsep}         % Abstand unten
  {\normalfont}     % Schrift des Bodys
  {0pt}             % Einschub der ersten Zeile
  {\bfseries}       % Darstellung von der Schrift in der Überschrift
  {:}               % Trennzeichen zwischen Überschrift und Body
  {.5em}            % Abstand nach dem Trennzeichen zum Body Text
  {\thmname{#3}}    % Name in eckigen Klammern
\theoremstyle{defi}

% ------ Definition zum Strich vor eines Texts
\newmdtheoremenv[
  hidealllines = true,       % Rahmen komplett ausblenden
  leftline = true,           % Linie links einschalten
  innertopmargin = 0pt,      % Abstand oben
  innerbottommargin = 4pt,   % Abstand unten
  innerrightmargin = 0pt,    % Abstand rechts
  linewidth = 3pt,           % Linienbreite
  linecolor = gray!40,       % Linienfarbe
]{defStrich}{Definition}     % Name der des formats "defStrich"

% ------ Definition zum Eck-Kasten um einen Text
\newmdtheoremenv[
  hidealllines = true,
  innertopmargin = 6pt,
  linecolor = gray!40,
  singleextra={              % Eck-Markierungen für die Definition
    \draw[line width=3pt,gray!50,line cap=rect] (O|-P) -- +(1cm,0pt);
    \draw[line width=3pt,gray!50,line cap=rect] (O|-P) -- +(0pt,-1cm);
    \draw[line width=3pt,gray!50,line cap=rect] (O-|P) -- +(-1cm,0pt);
    \draw[line width=3pt,gray!50,line cap=rect] (O-|P) -- +(0pt,1cm);
  }
]{defEckKasten}{Definition}  % Name der des formats "defEckKasten"  % Weitere Details sind ausgelagert

% ---- Für Todo Notes
\usepackage{todonotes}
\setlength {\marginparwidth }{2cm}      % Abstand für Todo Notizen

% ---- Zum Einbinden von PDF-Dokumenten
\usepackage{pdfpages}

% --- Zum Erstellen von Graphiken
\usepackage{tikz}
\usepackage{tabularx}
\usepackage{pgf-umlsd}
\usepackage{pgfplots}
\usepackage{adjustbox}
\usetikzlibrary{shadings}
\usetikzlibrary{backgrounds} 
\usetikzlibrary{fit, calc, matrix, shapes, positioning,patterns}
\pgfdeclarelayer{background}
\pgfsetlayers{background, main}
\pgfplotsset{compat = 1.9}
\definecolor{lightblue}{RGB}{173,216,230}
\definecolor{lightred}{RGB}{255,182,193}


% ---- Elektronische Version oder Gedruckte Version?
% ---- Unterschied: Die elektronische Version enthält keinen Platzhalter für die Unterschrift
\usepackage{ifthen}
\newboolean{e-Abgabe}
\setboolean{e-Abgabe}{false}    % false=gedruckte Fassung

% ---- Persönlichen Daten:
\newcommand{\titel}{Evaluationsmetriken von ML Algorithmen für Clustering, Klassifikation und Regression}
\newcommand{\titelheader}{}
\newcommand{\arbeit}{Ausarbeitung}
\newcommand{\studiengang}{Wirtschaftsinformatik}
\newcommand{\studienjahr}{2023}
\newcommand{\autor}{Jens Feser}
\newcommand{\autorReverse}{Feser, Jens}
\newcommand{\verfassungsort}{Mannheim}
\newcommand{\matrikelnr}{8658477}
\newcommand{\kurs}{WWI22SEB}
\newcommand{\bearbeitungsmonat}{}
\newcommand{\abgabe}{}
\newcommand{\bearbeitungszeitraum}{}
\newcommand{\firmaName}{SAP SE}
\newcommand{\firmaStrasse}{Dietmar-Hopp-Allee 16}
\newcommand{\firmaPlz}{69190 Walldorf, Deutschland}
\newcommand{\betreuerFirma}{}
\newcommand{\betreuerDhbw}{}

% ---- Metainformation für das PDF Dokument
\hypersetup{
	pdftitle    = {\titel},
	pdfsubject  = {\arbeit},
	pdfauthor   = {\autor},
	%pdfkeywords = {Keywords angeben},
	pdfcreator  = {LaTeX},
	%pdfproducer = {in der Regel pdfTeX}
}

% ---- Definition der Kopf- und Fußzeilen
\clearpairofpagestyles                          % Löschen von LaTeX Standard
\automark[section]{chapter}                     % Füllen von section und chapter
\renewcommand*{\chaptermarkformat}{}            % Entfernt die Kapitelnummer
\renewcommand*{\sectionmarkformat}{}            % Entfernt die Sectionnummer
% Angaben [für "plain"]{für "scrheadings"}
\ihead[]{\titelheader}                          % Kopfzeile links
\chead[]{}                                      % Kopfzeile mitte
\ohead[]{\rightmark}                            % Kopfzeile rechts
\ifoot[]{}                                      % Fußzeile links
\cfoot*{\sffamily\pagemark}                     % Fußzeile mitte
\ofoot[]{}                                      % Fußzeile rechts
\KOMAoptions{
   headsepline = 0.2pt,                         % Liniendicke Kopfzeile
   footsepline = false                          % Liniendicke Fußzeile
}
 % Kopf- und Fußzeile

% ---- Hilfreiches
\newcommand{\zB}{z.\,B. }   % "z.B." mit kleinem Leeraum dazwischen (ohne wäre nicht korrekt)
\newcommand{\dash}{d.\,h. }

\newcommand{\code}[1]{\texttt{#1}} % Ist einfacher zu schreiben als ständig \texttt und erlaubt
                                   % Änderungen im Nachhinein, wenn man z.B. Inline-Code anders stylen möchte.

% ---- Silbentrennung (falls LaTeX defaults falsch / nicht gewünscht sind)
\hyphenation{HANA}         % anstatt HA-NA
\hyphenation{Graph-Script} % anstatt GraphS-cript

% ---- Beginn des Dokuments
\begin{document}
\setlength{\parindent}{0pt}              % Keine Paragraphen Einrückung.
                                         % Dafür haben wir den Abstand zwischen den Paragraphen.
\setcounter{secnumdepth}{2}              % Nummerierungstiefe fürs Inhaltsverzeichnis
\setcounter{tocdepth}{2}                 % Tiefe des Inhaltsverzeichnisses. Ggf. so anpassen,
                                         % dass das Verzeichnis auf eine Seite passt.
\sffamily                                % Serifenlose Schrift verwenden.

% ---- Vorspann
% ------ Titelseite
\singlespacing
\thispagestyle{empty}
\begin{titlepage}
\enlargethispage{4cm}

\begin{figure}           % Logo vom Ausbildungsbetrieb und der DHBW
	% \vspace*{-5mm} % Sollte dein Titel zu lang werden, kannst du mit diesem "Hack" 
	%                  den Inhalt der Seite nach oben schieben.
	\begin{minipage}{0.49\textwidth}
		\flushleft
		\includegraphics[height=2.5cm]{Bilder/Logos/Logo_SAP.pdf} 
	\end{minipage}
	\hfill
	\begin{minipage}{0.49\textwidth}
		\flushright
		\includegraphics[height=2.5cm]{Bilder/Logos/Logo_DHBW.pdf} 
	\end{minipage}
\end{figure} 
\vspace*{0.1cm}

\begin{center}
	\huge{\textbf{\titel}}\\[1.5cm]
	\Large{\textbf{\arbeit}}\\[0.5cm]
	\normalsize{von}\\[1ex] \Large{\textbf{\autor}} \\[1cm]
\end{center}

\begin{center}
	\vfill
	\begin{tabular}{ll}
		Matrikelnummer, Kurs:            & \matrikelnr , \kurs \\[0.2cm]
	\end{tabular} 
\end{center}
\end{titlepage}
  % Titelseite
\newcounter{savepage}
\pagenumbering{Roman}                    % Römische Seitenzahlen
\onehalfspacing

% ------ Erklärung, Sperrvermerk, Abstact
\chapter*{Eidesstattliche Erklärung}
Ich versichere hiermit, dass ich meine \arbeit{} mit dem Thema:
\begin{quote}
	\textit{\titel}
\end{quote} 
gemäß § 5 der \enquote{Studien- und Prüfungsordnung DHBW Technik} vom 29. September 2017 selbstständig verfasst und keine anderen als die angegebenen Quellen und Hilfsmittel benutzt habe. Die Arbeit wurde bisher keiner anderen Prüfungsbehörde vorgelegt und auch nicht veröffentlicht.


\vspace{1cm}

\verfassungsort, den \today \\[0.5cm]
\ifthenelse{\boolean{e-Abgabe}}
	{\underline{Gez. \autor}}
	{\makebox[6cm]{\hrulefill}}\\ 
\autorReverse

\chapter*{Disclaimer}

Der Umfang dieser Arbeit von 14 Seiten ergibt sich aus der Vielzahl an umfangreichen mathematischen LaTeX-Darstellungen sowie den zahlreichen Beispielen, die zur Veranschaulichung der behandelten Konzepte dienen.

Der Quellcode für das Projekt sowie die zugehörige Demo sind auf GitHub zu finden:  

\url{https://github.com/Jens011203/Evaluierung_Integrationsseminar}


% ------ Inhaltsverzeichnis
\singlespacing
\tableofcontents

% ------ Verzeichnisse
\renewcommand*{\chapterpagestyle}{plain} 
\pagestyle{plain}
\chapter*{Abkürzungsverzeichnis}
\addcontentsline{toc}{chapter}{Abkürzungsverzeichnis} % Hinzufügen zum Inhaltsverzeichnis 

\begin{acronym}[WYSISWG] % längstes Kürzel wird verw. für den Abstand zw. Kürzel u. Text

	% Alphabetisch selbst sortieren - nicht verwendete Kürzel rausnehmen!
  % \acro{TEST}{Test Ausgeschrieben} Vorlage Abkürzung
  \acro{DBI}{Davies-Bouldin-Index}
  \acro{ARI}{Adjusted Rand Index}
  \acro{AMI}{Adjusted Mutal Information}
  \acro{MI}{Mutal Information}
  \acro{TP}{True Positive}
  \acro{FP}{False Positive}
  \acro{TN}{True Negative}
  \acro{FN}{False Negative}
  \acro{ROC}{Receiver Operating Characteristic}
  \acro{AUC}{Area under the Curve}
  \acro{MSE}{Mean Squared Error}
  \acro{RMSE}{Root Mean Squared Error}
  \acro{MAE}{Mean Absolute Error}

\end{acronym} 

\listoffigures                          % Erzeugen des Abbildungsverzeichnisses 
\setcounter{savepage}{\value{page}}


% ---- Inhalt der Arbeit
\cleardoublepage
\pagenumbering{arabic}                  % Arabische Seitenzahlen für den Hauptteil
\setlength{\parskip}{0.5\baselineskip}  % Abstand zwischen Absätzen
\rmfamily
\renewcommand*{\chapterpagestyle}{scrheadings}
\pagestyle{scrheadings}
\onehalfspacing
\chapter{Einleitung}

Die Evaluierung von Machine-Learning-Modellen erfolgt anhand geeigneter
Metriken. Dabei hängt die Wahl der Metriken vom
spezifischen Aufgabentyp ab (Clustering, Klassifikation oder Regression). Gute Metriken
legen fest, wie Qualität im jeweiligen Kontext gemessen wird. So strebt man
beim Clustering dichte, gut getrennte Cluster an, bei der Klassifikation eine
hohe Übereinstimmung mit den Klassenlabels, und bei der Regression geringe
Abweichungen der Vorhersagen vom wahren Wert. \cite{Miller2024} legt dar,
dass das Verständnis und die richtige Auswahl von Metriken entscheidend sind,
um die Modellleistung objektiv zu beurteilen und Fehlinterpretationen zu
vermeiden. 

Die vorliegende Arbeit soll einen Überblick über die wichtigsten Evaluationsmetriken in diesen Zusammenhang liefern.
Sie soll dabei als Grundlage dienen, welche Evaluationsmetriken für die verschiedenen ML-Algorithmen verwendet werden können.

Im Folgenden werden die wichtigsten Metriken für Clustering, Klassifikation und Regression vorgestellt
und hinsichtlich ihrer Stärken und Schwächen diskutiert.


\chapter{Clustering}
Clustering bezeichnet Algorithmen im Bereich des \emph{unsupervised Learnings},
welche Datenpunkte anhand ihrer Ähnlichkeit zueinander in Gruppen bzw. Cluster
ordnen.
Gute Cluster zeichnen sich durch eine hohe Kohäsion innerhalb
eines Clusters und gute Separation zwischen den Clustern aus. Typische
Evaluationsziele sind daher, kompakte, dicht gepackte Cluster zu bilden, die
gleichzeitig weit voneinander getrennt sind (geringe Intra-Cluster-Distanz,
hohe Inter-Cluster-Distanz). 
Bei der Evaluierung von Clustering-Algorithmen gibt es zwei Möglichkeiten (vgl. \cite{Miller2024}):
\begin{itemize}
 \item \textbf{extrinsische Metriken:} Ist eine sogenannte \textit{Ground Truth} (also die wahre Zuordnung der Daten in Cluster) verfügbar, kann diese zu Evaluierung
    der Algorithmen verwendet werden. Extrinsische Metriken erzielen gute Werte, wenn die vorhergesagten Cluster sehr ähnlich zu den tatsächlichen Clustern sind.
 \item \textbf{intrinsische Metriken:} Diese werden verwendet, wenn keine \textit{Ground Truth} in den Daten vorhanden ist. Zur Bewertung der Algorithmen wird 
    die Ähnlichkeit von Punkten im selben Cluster mit der Ähnlichkeit zu anderen Clustern verglichen. Intrinsische Metriken liefern gut Werte, wenn die 
    Intra-Cluster-Ähnlichkeit größer ist als Inter-Cluster-Ähnlichkeit
\end{itemize}
Das Beispiel mit komplexen 
halbmondförmigen Clustern zeigt die Schwäche von intrinsischen Metriken auf. Die intrinsischen Metriken bewerten den Clustering-Algorithmus
der die Cluster möglichst kugelförmig vorhersagt besser, obwohl der andere Algorithmus die Cluster perfekt vorhersagt. Deshalb wenn eine 
\emph{Ground Truth} vorhanden ist, sollten eher extrinsische Metriken verwendet werden.

\url{https://github.com/Jens011203/Evaluierung_Integrationsseminar/blob/main/src/Clustering.ipynb}

Im Folgenden werden relevante Metriken vorgestellt:

\section{Intrinsische Metriken}
\subsection{Silhouette-Koeffizient} Der Silhouette-Koeffizient bewertet für jeden
Datenpunkt $i$ dessen Zugehörigkeit zu einem Cluster. Sei $a(i)$ der
durchschnittliche Abstand des Objekts $i$ zu allen anderen Punkten desselben
Clusters, und $b(i)$ der kleinste durchschnittliche Abstand zu Punkten zum nächstgelegenen Fremdcluster.
Dann ist der Silhouette-Wert definiert als 
\[ s(i) =
\frac{b(i) - a(i)}{\max\{a(i),b(i)\}}, \] 

wobei $s(i)\in[-1,1]$. Werte
nahe 1 bedeuten, dass $i$ gut zum eigenen Cluster passt und weit von anderen
Clustern entfernt ist, Werte nahe $0$ deuten auf überlappende Cluster hin und
negative Werte darauf, dass $i$ eher falsch zugeordnet.
Der Gesamtsilhouette-Index wird als Mittelwert $\overline{s} =
\frac{1}{N}\sum_i s(i)$ über alle Objekte berechnet (vgl. \cite{Rousseeuw1987} S.55f). 

\paragraph{Beispiel:} Betrachten wir eine stark vereinfachte Situation mit zwei Clustern A und B (siehe \ref{fig:silhouette_example}):
\begin{itemize}
  \item Cluster A enthält die beiden Punkte $x_1$ und $x_2$.
  \item Cluster B enthält die beiden Punkte $y_1$ und $y_2$.
  \item Punkt $i$ ist aktuell Cluster A zugeordnet.
\end{itemize}
Die Abstände betragen:
\[
  d_1 = \|i - x_1\| = 1.3,
  \quad d_2 = \|i - x_2\| = 0.7,
  \quad d_3 = \|i - y_1\| = 2.8,
  \quad d_4 = \|i - y_2\| = 3.2.
\]
Somit ergeben sich folgende mittlere Abstände:
\[
  a(i) = \frac{d_1 + d_2}{2} = \frac{1.3 + 0.7}{2} = 1
  b(i) = \frac{d_3 + d_4}{2} = \frac{2,8 + 3.2}{2} = 3.
\]
Somit ergibt sich
\[
  s(i) = \frac{b(i) - a(i)}{\max\{a(i), b(i)\}} = \frac{3 - 1}{3} \approx 0.66.
\]
Ein Wert $s(i) \approx 0.66$ zeigt, dass $i$ relativ gut in Cluster A integriert ist, da $a(i) \ll b(i)$.

\begin{figure}[ht]
  \centering
  \begin{tikzpicture}[scale=1.5, every node/.style={font=\small}]
    % Cluster A (blau): zwei Punkte
    \fill[blue!60] (0,0) circle (3pt) node[below left] {$x_1$};
    \fill[blue!60] (0.5,0.2) circle (3pt) node[above left] {$x_2$};
    % Cluster B (rot): ein Punkt
    \fill[red!60] (3,0) circle (3pt) node[below right] {$y_1$};
    \fill[red!60] (3.2,0.3) circle (3pt) node[above right] {$y_2$};
    % Punkt i (grün)
    \coordinate (i) at (1,0);
    \fill[green!70!black] (i) circle (3pt) node[above] {$i$};

    % a(i)-Linien zu beiden A-Punkten
    \draw[blue!60, dashed, <->] (i) -- node[below] {} (0,0);
    \draw[blue!60, dashed, <->] (i) -- node[above] {} (0.5,0.2);

    % b(i)-Linie
    \draw[red!60, dashed, <->] (i) -- node[below] {} (3,0);
    \draw[red!60, dashed, <->] (i) -- node[above] {} (3.2,0.3);

    % Legende
    \begin{scope}[shift={(0,-1)}]
      \fill[blue!60] (0,0) circle (3pt) node[right=4pt] {Cluster A};
      \fill[red!60] (2,0) circle (3pt) node[right=4pt] {Cluster B};
      \fill[green!70!black] (4,0) circle (3pt) node[right=4pt] {Punkt $i$};
    \end{scope}
  \end{tikzpicture}
  \caption{Beispiel einer einfachen Clusterzuordnung}
  \label{fig:silhouette_example}
\end{figure}


\paragraph{Kritische Einordnung:}
Ein Nachteil des Silhouette‐Index ist, dass er runde und gleichmäßig dichte Cluster und runde
Cluster annimmt. Bei stark ungleich geformten oder  
verrauschten Clustern kann er daher irreführende Werte liefern (Siehe komplexeres Code-beispiel im verlinkten Beispiel mit halbmondförmigen Cluster: \url{https://github.com/Jens011203/Evaluierung_Integrationsseminar/blob/main/src/Clustering.ipynb}).
Außerdem reagiert er nicht auf Merkmals‑Bias: Trennt ein Algorithmus auf Basis irrelevanter Variablen,  
bleibt dies im Silhouette‑Score unentdeckt (vgl. \cite{Miller2024}). Deswegen wird diese Metrik vor allem genutzt wenn keine wahren labels in den Daten vorhanden sind.

\subsection{Davies-Bouldin-Index} Der \ac{DBI} misst
die Ähnlichkeit zwischen jedem
Cluster und dem ähnlichsten Fremdcluster. Für jeden Cluster $i$ definiert
man $s_i$ als den durchschnittlichen Abstand der Punkte von ihrem
Cluster-Zentroid (also die Varianz im Cluster), und für Clusterpaare $(i,j)$ den Abstand $d_{ij}$ zwischen
den Zentroiden. Dann berechnet man für jedes Paar $R_{ij} = (s_i + s_j)/d_{ij}$
und wählt für jeden Cluster $i$ den größten Wert $\max_{j\neq i} R_{ij}$. Der
Davies-Bouldin-Index ist das Mittel dieser Maxima: 
\[ DBI = \frac{1}{k}\sum_{i=1}^k \max_{j\neq i}\frac{s_i + s_j}{d_{ij}}, \]
mit $k$ der Anzahl der Cluster. Ein niedriger \ac{DBI} bedeutet bessere Cluster (maximale
Separation und geringe Streuung). Der Wert kann theoretisch bis $0$ reichen
(perfekte Trennung) (vgl. \cite{Davies1979} S.224f).

\paragraph{Beispiel:} Nehmen wir zwei Cluster mit Zentroiden \(c_1\) und \(c_2\),
deren Streuungen \(s_1=1\) und \(s_2=0{,}5\) sind, und Abstand \(d_{12}=2\). Dann ist
\[
R_{12}=\frac{s_1 + s_2}{d_{12}} = \frac{1 + 0.5}{2} = 0.75,
\]
was für dieses Paar den hohen Ähnlichkeitswert beiträgt und den DBI erhöht.
Befänden sich die Zentroiden jedoch bei \(d_{12}=5\), so wäre  
\[
R_{12} = \frac{1 + 0.5}{5} = 0.3,
\]
und der \ac{DBI} bliebe gering (gutes Clustering).

\paragraph{Kritische Einordnung:} Der \ac{DBI} ist ein effizientes, intrinsisches Maß, bei dem kleinere Werte
eine bessere Clusterstruktur anzeigen (Minimum = 0). Im Vergleich zum Silhouette‐Index ist die Berechnung
deutlich schneller und daher für großvolumige Datensätze sehr gut geeignet.
Nachteil: Er setzt euklidische Abstände zwischen Clusterzentroiden voraus und kann bei spärlichen oder
nicht‑euklidischen Daten irreführende Ergebnisse liefern, wenn alternative Distanzmaße besser zur Struktur der Daten passen.
(vgl. \cite{Miller2024}). Zudem kann er wie der Silhoutte-Koeffizient bei komplexeren Cluster-Strukturen verfälschte Werte liefern.

\subsection{Calinski-Harabasz-Index} Der Calinski-Harabasz-Index (auch Variance
Ratio Criterion) setzt die Streuung zwischen den Clustern ins Verhältnis zur
Streuung innerhalb der Cluster (vgl. \cite{Calinski1974}). Formal
sei $N$ die Gesamtzahl der Punkte, $k$ die Anzahl der Cluster, $B$ die
zwischen-Cluster-Varianz (Summe der quadrierten Abstände der Cluster-Zentren
zum Gesamtdurchschnitt) und $W$ die intra-Cluster-Varianz (Summe der
quadrierten Abstände der Punkte zu ihren jeweiligen Zentren). Dann ist der
CH-Index definiert als 
\[
CH 
= \frac{B/(k-1)}{W/(N-k)}
= \frac{B}{W} \cdot \frac{N - k}{k - 1}
\]
Ein höherer Wert
zeigt klarere Trennung und dichtere Cluster gute Klassifikation, während ein
kleinerer Wert auf schlechtere Partitionen (vgl. \cite{Calinski1974}). 

\paragraph{Beispiel:} Sind die Cluster perfekt kompakt ($W\to0$), wächst $CH$ ins
Unendliche. Bei nur einem Cluster ($k=1$) definiert man meist $CH=0$.
Üblicherweise wählt man das $k$, für das $CH$ am größten ist.

\paragraph{Kritische Einordnung:} Der CH-Index ist schnell berechenbar und eignet
sich für große Datensätze. Wie auch die anderen intrinsischen Metriken belohnt er kompakte, klar getrennte Cluster, kann
aber bei sehr vielen kleinen Clustern hohe Werten liefern und ist nur im
Vergleich unterschiedlicher $k$-Lösungen sinnvoll (vgl. \cite{Miller2024}).

\section{Extrinsische Metriken}

\subsection{Purity}
Die Purity eines Clusters \(C_i\) mit \(n_i\) Datenpunkten und
\(n_{h}^{(i)}\) Datenpunkten der dominanten (wahrheitsgemäßen) Klasse ist
definiert als {vgl. \cite{huang2008similarity} S.54}:
\[
  P(C_i)=\frac{1}{n_i}\max_h n_{h}^{(i)}.
\]
Der Gesamt‑Purity-Score eines Clusterings ist das gewichtete Mittel über alle \(k\) Cluster:
\[
  \text{Purity} = \sum_{i=1}^k \frac{n_i}{N}\,P(C_i).
\]

\paragraph{Beispiel:}  
Ein Vorgeschlagenes Cluster \(C_1\) enthält 10 Punkte, davon 7 aus Klasse A.
Dann gilt
\[
  P(C_1) = \frac{7}{10} = 0{,}7.
\]

\paragraph{Kritische Einordnung:}  
Purity ist sehr einfach und anschaulich, aber durch Einpunkt‑Cluster trivialisierbar und ignoriert Minderheitsklassen.
Zudem ist ein fairer Vergleich nur bei gleicher Anzahl \(k\) möglich.

\subsection{Adjusted Rand Index (ARI)}
Der \ac{ARI} beschreibt die Übereinstimmung zweier Cluster \(X\) (Vorhersage),
\(Y\) (\emph{Ground Truth}) und korrigiert um die erwartete Übereinstimmung
bei zufälliger Zuordnung (vgl. \cite{Hubert1985} S.194ff und \cite{Miller2024}):
\[
\mathrm{ARI}
= \frac{\mathrm{RI} - \mathbb{E}[\mathrm{RI}]}{1 - \mathbb{E}[\mathrm{RI}]},
\]
wobei der \emph{Rand‐Index}
\[
\mathrm{RI} = \frac{a + d}{\binom{n}{2}}
\]
den Anteil aller Paare misst, die in beiden Partitionen entweder im gleichen Cluster (\(a\)) oder in verschiedenen Clustern (\(d\)) liegen.  
Die Erwartung \(\mathbb{E}[\mathrm{RI}]\) unter zufälliger Clusterbildung berechnet sich aus den Clustergrößen in \(X\) und \(Y\):
\[
\mathbb{E}[\mathrm{RI}]
= \frac{\displaystyle\sum_i \binom{n_i^{(X)}}{2}\;\times\;\sum_j\binom{n_j^{(Y)}}{2}}{\binom{n}{2}^2},
\]
wobei \(n_i^{(X)}\) und \(n_j^{(Y)}\) die Größen der Cluster \(i\) in \(X\) bzw.\ \(j\) in \(Y\) sind.

\textbf{Wertebereich:} \(\mathrm{ARI}\in[-1,1]\), wobei  
\(\mathrm{ARI}=1\) perfekte Übereinstimmung,  
\(\mathrm{ARI}=0\) Zufallsniveau und  
\(\mathrm{ARI}<0\) schlechter als Zufall bedeutet.

\paragraph{Beispiel:}

Gegeben sind drei Objekte \(\{1,2,3\}\) mit  
\[
X:\{1,2\},\{3\}, 
\quad
Y:\{1\},\{2,3\}.
\]
Es existieren \(\binom{3}{2}=3\) Paare. Wir ermitteln
\[
a=0\quad(\text{kein Paar in beiden Partitionen im selben Cluster}),
\]
\[
d=1\quad(\text{z.\,B. Paar }(1,3)\text{ in beiden Partitionen getrennt}).
\]
Damit ist:
\[
\mathrm{RI}=\frac{a+d}{3}=\frac{0+1}{3}=\tfrac{1}{3}.
\]
Für die Erwartung unter Zufall:
\[
\sum_i\binom{n_i^{(X)}}{2}=1,\quad
\sum_j\binom{n_j^{(Y)}}{2}=1,
\quad
\mathbb{E}[\mathrm{RI}]=\frac{1\cdot1}{3^2}=\tfrac{1}{9}.
\]
Folglich
\[
\mathrm{ARI}
=\frac{\tfrac13-\tfrac19}{1-\tfrac19}
=\frac{\tfrac{2}{9}}{\tfrac{8}{9}}
=\tfrac{1}{4}
=0{,}25.
\]

\paragraph{Kritische Einordnung:}  
\ac{ARI} ist eine sehr häufig verwendete und zuverlässige Metrik. Wie bei allen extrinsischen Metriken ist hier zur Bewertung eine \textit{Ground Truth}
notwendig. Eine weitere Einschränkung beim \ac{ARI} ist die Verzerrung bezüglich der Clustergröße. 
Wenn eine Clustering-Lösung eine Mischung aus großen und kleinen Clustern enthält, wird \ac{ARI} vorwiegend von den großen Clustern beeinflusst (vgl. \cite{Miller2024}).

\subsection{Adjusted Mutual Information (AMI)}
Die \ac{AMI} beschreibt wie viel Information zwischen der dem vorhergesagten
Cluster \(X\) und den tatsächlichen Cluster Y geteilt wird und korrigiert für
zufällige Übereinstimmungen. 
Die Metrik hat denselben Wertebereich wie \ac{ARI} und wird wie folgt bestimmt (vgl. \cite{Miller2024}):
\[
\mathrm{AMI}(X,Y) = \frac{\mathrm{MI}(X,Y) - \mathbb{E}[\mathrm{MI}(X,Y)]}{\mathrm{avg}(H(X),H(Y)) - \mathbb{E}[\mathrm{MI}(X,Y)]},
\]
wobei:
\begin{itemize}
    \item $H$ = individuelle Entropie -- ein Maß für die erwartete Unsicherheit
    \item $\mathrm{MI}$ = Der Mutual-Information-Algorithmus 
    \item $\mathbb{E}$ = der Erwartungswert basierend auf Zufall
\end{itemize}

\paragraph{Beispiel:}
Gegeben sind die Clusterings \(X: \{A,B\}, \{C\}\) und \(Y: \{A\}, \{B,C\}\).
Würde man hiervon die \ac{MI} berechnen erhält man \(MI = 0{,}251\). 
Nach Korrektur um den Zufallswert erhält man \(\mathrm{AMI}(X,Y) = 0{,}136\), was etwas besser als 0 (zufällig) ist.

\paragraph{Kritische Einordnung:}
\ac{AMI} bietet durch die Zufallskorrektur eine robustere Interpretation als
die einfache \ac{MI}. Ein wesentlicher Unterschied zu \ac{ARI} liegt in den
jeweiligen Verzerrungen: Während \ac{ARI} Lösungen mit ähnlich großen Clustern
bevorzugt, ist \ac{AMI} zu „reinen“ Clustern verzerrt, die nur einen Klassentyp
enthalten und oft unausgewogen sind (vgl. \cite{Miller2024}). 

\chapter{Klassifikation}

Klassifikation bezeichnet Algorithmen im Bereich des \emph{supervised
Learnings}, welche Datenpunkte anhand bekannter Klassenlabels in vordefinierte
Kategorien einordnen. Ziel ist, auf Basis eines Trainingsdatensatzes mit
Features und zugehörigen Labels ein Modell zu lernen, das für neue, unbekannte
Datenpunkte möglichst korrekte Vorhersagen trifft. Gute Klassifikatoren
zeichnen sich durch hohe Treffergenauigkeit (Accuracy), gute Trennung
zwischen den Klassen (Precision/Recall, AUC) und robuste Leistung bei
unterschiedlichen Klassenhäufigkeiten aus (vgl \cite{Miller2024}).

Viele dieser Metriken basieren im binären Fall auf den vier Kategorien einer
\emph{Konfusionsmatrix} (\ac{TP}, \ac{FP}, \ac{TN}, \ac{FN})
(vgl. Abbildung~\ref{fig:confmat}). Im folgenden Beispiel wurden von 100 Instanzen
insgesamt 30 als \textbf{True Positives}, 10 als \textbf{False Positives}, 15 als \textbf{False Negatives}
und 45 als \textbf{True Negatives} klassifiziert. Dieses Beispiel wird im Folgenden zur als Beispiel
der Berechnung einzelner Metriken genutzt.
Ein komplexeres Beispiel, in welchen Algorithmen auf einen Brustkrebsdatensatz trainiert wurden und bestimmen sollen ob der 
Brustkrebs bösartig oder gutartig ist, ist im verlinkten GitHub-Repository zu finden.

\begin{figure}[ht]
  \centering
  \begin{tikzpicture}[font=\small]
    % Gitter und Linien
    \draw (0,0) rectangle (4,4);
    \draw (2,0) -- (2,4); % vertikale Mittellinie
    \draw (0,2) -- (4,2); % horizontale Mittellinie

    % Boxenbeschriftung + Beispielwerte
    \node at (1,3.2) {\textbf{TP}};
    \node at (1,2.6) {(=30)};
    
    \node at (3,3.2) {\textbf{FP}};
    \node at (3,2.6) {(=10)};
    
    \node at (1,1.2) {\textbf{FN}};
    \node at (1,0.6) {(=15)};
    
    \node at (3,1.2) {\textbf{TN}};
    \node at (3,0.6) {(=45)};

    % Achsenbeschriftung
    \node[rotate=90] at (5,2) {Tatsächliche Klasse};
    \node at (2,-1) {Vorhergesagte Klasse};

    % Spaltenbeschriftung (Predicted)
    \node at (1,4.3) {Positiv};
    \node at (3,4.3) {Negativ};

    % Zeilenbeschriftung (Actual)
    \node at (-0.8,3) {Positiv};
    \node at (-0.8,1) {Negativ};
  \end{tikzpicture}
  \caption{Konfusionsmatrix mit Beispielwerten}
  \label{fig:confmat}
\end{figure}

\section{Accuracy (Treffergenauigkeit)}
Die Accuracy misst den Anteil korrekt klassifizierter Instanzen an allen Fällen.  
\[
\mathrm{Accuracy}=\frac{TP+TN}{TP+FP+FN+TN} = \frac{30+45}{100} = 0{,}75
\]
Somit misst sie wie viele Vorhersagen insgesamt richtig waren.
Eine hohe Accuracy ist nur aussagekräftig, wenn die Klassenverteilung ausgeglichen ist (vgl. \cite{Miller2024}).

\section{Precision, Recall, F1-Score}
Diese drei Metriken sind besonders wichtig, wenn die Klassenverteilung unausgeglichen ist (z.\,B. bei seltenen Ereignissen)(vgl. \cite{powers2020evaluation} S.38f).

\textbf{Precision} (Positiver Vorhersagewert):
\[
\mathrm{Precision} = \frac{TP}{TP+FP} = \frac{30}{30+10} = 0{,}75
\]
Precision beschreibt den Anteil der als positiv vorhergesagten Fälle, die tatsächlich positiv sind.  
Somit sollte diese Metrik verwendet werden, wenn \ac{FP} teuer sind (z.\,B. Spam‑Filter).  
Allerdings wird ignoriert, wie viele echte Positive verpasst werden.

\textbf{Recall} (Sensitivität, Trefferquote):
\[
\mathrm{Recall} = \frac{TP}{TP+FN} = \frac{30}{30+15} = 0{,}67
\]
Recall beschreibt den Anteil der tatsächlichen Positiven, die korrekt erkannt wurden.  
Die Metrik ist somit relevant, wenn \ac{FN} kritisch sind (z.\,B. Krebsdiagnose).  

\textbf{F\(_1\)-Score} (harmonisches Mittel):
\[
F_1 = 2 \cdot \frac{\mathrm{Precision}\,\times\,\mathrm{Recall}}{\mathrm{Precision} + \mathrm{Recall}}
\approx 0{,}71
\]
Der F1-Score balanciert Precision und Recall und gibt nur dann hohe Werte, wenn beide hoch sind.  
Beachtet allerdings \ac{TN} nicht und ist daher nicht vollständig aussagekräftig, wenn TN wichtig sind (vgl. \cite{powers2020evaluation} S38f).

\section{ROC-Kurve und AUC (Area Under the Curve)}
Die \ac{ROC}-Kurve ist ein Werkzeug zur Bewertung von binären
Klassifikationsmodellen. Sie zeigt den Zusammenhang zwischen der \emph{True
Positive Rate} (Recall) und der \emph{False Positive Rate} (FPR) für
verschiedene Entscheidungs­schwellen (Thresholds) des Modells.

\begin{align*}
\mathrm{TPR} &= \frac{TP}{TP + FN} \quad \text{(Recall)} \\
\mathrm{FPR} &= \frac{FP}{FP + TN}
\end{align*}

\textbf{Interpretation:}  
Die \ac{ROC}-Kurve zeigt, wie gut ein Modell zwischen den Klassen unterscheidet. Sie
beginnt immer bei (0,0) und endet bei (1,1). Je näher die Kurve an der oberen
linken Ecke liegt, desto besser ist das Modell (vgl. \cite{Miller2024}).

\ac{AUC} ist die Fläche unter der ROC-Kurve:
\begin{itemize}
  \item AUC = 1: perfekter Klassifikator
  \item AUC = 0.5: reines Raten (Zufall)
  \item AUC < 0.5: systematisch falsche Klassifikation
\end{itemize}

\textbf{Vorteile:}  
AUC ist unabhängig vom gewählten Schwellenwert und eignet sich besonders bei unausgeglichenen Klassenverhältnissen.

\begin{figure}[ht]
  \centering
  \begin{tikzpicture}[scale=1.1, font=\small]
    % Achsen
    \draw[->] (0,0) -- (4,0) node[right] {False Positive Rate (FPR)};
    \draw[->] (0,0) -- (0,4) node[above] {True Positive Rate (TPR)};

    % Diagonale
    \draw[dashed] (0,0) -- (3.5,3.5);

    % Beispiel ROC-Kurve
    \draw[very thick, blue] 
      (0,0) -- (0,1.5) -- (1,2.5) -- (2,3) -- (3.5,3.5);

    % Fläche unter Kurve (schematisch)
    \fill[blue!10] 
      (0,0) -- (0,1.5) -- (1,2.5) -- (2,3) -- (3.5,3.5) -- (0,0);

    % Markierungen
    \foreach \i/\label in {0/0, 1/0.25, 2/0.5, 3/0.75}
      {
        \draw (\i,0) -- (\i,-0.1) node[below] {\label};
        \draw (0,\i) -- (-0.1,\i) node[left] {\label};
      }
    \draw (3.5,0) -- (3.5,-0.1) node[below] {1};
    \draw (0,3.5) -- (-0.1,3.5) node[left] {1};

  \end{tikzpicture}
  \caption{Beispielhafte ROC-Kurve (blau) und zufälliger Klassifikator (gestrichelt)}
  \label{fig:roc}
\end{figure}

\textbf{Beispiel:}  
Angenommen, ein Modell trifft mit variabler Schwelle folgende Entscheidungen (sortiert nach Modell-Score):

\begin{center}
\begin{tabular}{|c|c|c|}
\hline
Score & Wahre Klasse & Klassifikation bei Schwelle \\
\hline
0.9 & Positiv & Positiv \\
0.8 & Negativ & Positiv \\
0.7 & Positiv & Positiv \\
0.4 & Positiv & Negativ \\
0.2 & Negativ & Negativ \\
\hline
\end{tabular}
\end{center}

Bei Variation der Schwelle kann man verschiedene (FPR, TPR)-Punkte erzeugen und so die ROC-Kurve aufbauen (siehe Abbildung \ref{fig:roc}).

\section{Mehrklassige Klassifikation}

Zur Evaluation nicht‑binärer Klassifikationsprobleme lassen sich die Evaluationsmetriken über alle Klassen hinweg zusammenfassen.
Hierbei gibt es zwei Ansätze (vgl. \cite{takahashi2022confidence}):

\begin{itemize}
  \item \textbf{Micro-Averaging}:  
    \begin{itemize}
      \item Man fasst alle Klassen‑Confusion‑Matrix‑Einträge instanzübergreifend zusammen:
      \[
        TP = \sum_{i=1}^r TP_i,\quad
        FP = \sum_{i=1}^r FP_i,\quad
        FN = \sum_{i=1}^r FN_i.
      \]
      \item Anschließend werden die Metriken wie im binären Fall berechnet.
    \end{itemize}

  \item \textbf{Macro-Averaging}:  
    \begin{itemize}
      \item Für jede Klasse \(i\) separat:
      \[
        P_i = \frac{TP_i}{TP_i + FP_i},\quad
        R_i = \frac{TP_i}{TP_i + FN_i},\quad
      \]
      \[
        F_{1,i} = 2\,\frac{P_i\,R_i}{P_i + R_i}.
      \]
      \item Die \emph{Macro-F\(_1\)} ist der einfache Mittelwert:
      \[
        \mathrm{F}_{1,\text{macro}}
        = \frac{1}{r} \sum_{i=1}^r F_{1,i}.
      \]
    \end{itemize}
\end{itemize}

Zusammengefasst gilt für die beiden Metriken:
\begin{itemize}
  \item \textbf{Micro-F\(_1\)}: Gesamtüberblick, instanzgewichtet, robust bei Klassenungleichgewicht.
  \item \textbf{Macro-F\(_1\)}: Klassen gleichgewichtet, macht Schwächen seltener Klassen sichtbar.
\end{itemize}

\chapter{Regression}
Regression bezeichnet Algorithmen im Bereich des \emph{supervised Learnings},
welche kardinal skalierte Zielwerte \(y\) vorhersagen. Ziel
ist, auf Basis eines Trainingsdatensatzes mit Eingabevektoren \(\mathbf{x}_i\)
und zugehörigen Messwerten \(y_i\) ein Modell zu erlernen, das für neue
Datenpunkte möglichst geringe Abweichungen \(\hat y_i - y_i\) erzielt. Anders
als bei der Klassifikation interessiert hier nicht die korrekte
Kategorisierung, sondern die Minimierung des Vorhersagefehlers (vgl. \cite{Miller2024}). 
Ein komplexeres Beispiel im verlinkten GitHub-Repository trainiert und evaluiert zwei einfachere Regression-Algorithmen 
anhand eines Diabetes-Datensatzes.

Im Folgenden werden die gebräuchlichsten Fehlermaße vorgestellt.

\section{\acs{MSE} und \acs{RMSE}}

Der \ac{MSE} ist definiert als 
\[
  \mathrm{MSE} = \frac{1}{N}\sum_{i=1}^N (\hat y_i - y_i)^2,
\]
der \ac{RMSE} als
\[
  \mathrm{RMSE} = \sqrt{\mathrm{MSE}}.
\]
Beide Metriken messen den durchschnittlichen quadratischen Fehler und sind in der Einheit von \(y\) (\ac{RMSE}) bzw. im Quadrat davon (\ac{MSE}).
Folglich deuten niedrigere Werte auf eine bessere Regression hin.

\paragraph{Beispiel:}  
Angenommen, wir haben \(N=3\) Testwerte mit  
\((y_1,y_2,y_3) = (2,\,5,\,7)\) und Vorhersagen \((\hat y_1,\hat y_2,\hat y_3) = (3,\,4,\,10)\).  
Dann rechnen wir:
\[
  (\hat y_1 - y_1)^2 = (3-2)^2 = 1,\quad
  (\hat y_2 - y_2)^2 = (4-5)^2 = 1,\quad
  (\hat y_3 - y_3)^2 = (10-7)^2 = 9.
\]
Somit ist
\[
  \mathrm{MSE} = \frac{1 + 1 + 9}{3} = \frac{11}{3} \approx 3{,}67,
  \quad
  \mathrm{RMSE} = \sqrt{3{,}67} \approx 1{,}92.
\]
Somit liegt der Fehler der Vorhersagen bei etwa 1,92 Einheiten.

\paragraph{Kritische Einordnung}
\ac{MSE}/\ac{RMSE} bestrafen größere Fehler stärker (Quadrierung) und sind leicht zu optimieren, reagieren jedoch empfindlich auf Ausreißer (vgl. \cite{hodson2022root}).

\section{\ac{MAE}}

Der \ac{MAE} ist definiert als
\[
  \mathrm{MAE} = \frac{1}{N}\sum_{i=1}^N \bigl|\hat y_i - y_i\bigr|.
\]
Er misst den durchschnittlichen absoluten Fehler in derselben Einheit wie \(y\). Somit deuten auch hier niedrigere Werte auf eine bessere Regression hin.

\paragraph{Beispiel:}  
Mit den gleichen Werten \((y_i) = (2,5,7)\) und \((\hat y_i) = (3,4,10)\) erhalten wir:
\[
  |3-2| = 1,\quad |4-5| = 1,\quad |10-7| = 3,
\]
\[
  \mathrm{MAE} = \frac{1 + 1 + 3}{3} = \frac{5}{3} \approx 1{,}67.
\]
Somit weichen die Vorhersagen um etwa \(1{,}67\) von den Vorhersagen ab.

\paragraph{Kritische Einordnung}
\ac{MAE} ist einfacher zu interpretieren und robuster gegenüber Ausreißern, gewichtet aber alle Fehler linear (vgl. \cite{hodson2022root}).

\section{\(R^2\) (Bestimmtheitsmaß)}

Das Bestimmtheitsmaß \(R^2\) gibt an, welcher Anteil der Varianz der Zielwerte durch das Modell erklärt wird:
\[
  R^2 = 1 - \frac{\sum_{i=1}^N (\hat y_i - y_i)^2}{\sum_{i=1}^N (y_i - \bar y)^2},
  \quad
  \bar y = \frac{1}{N}\sum_{i=1}^N y_i.
\]
Die Werte liegen zwischen \(-\infty\) und \(1\).\\

\paragraph{Beispiel:}  
Für unsere Werte \((y_i) = (2,5,7)\) ist der Mittelwert 
\(\bar y = \tfrac{2+5+7}{3} = 14/3 \approx 4{,}67.\)  
Berechne die Quadratsummen:
\[
  \sum(\hat y_i - y_i)^2 = 1 + 1 + 9 = 11,
  \quad
  \sum(y_i - \bar y)^2 
    = (2-4{,}67)^2 + (5-4{,}67)^2 + (7-4{,}67)^2
    \approx 7{,}11 + 0{,}11 + 5{,}44 = 12{,}67.
\]
Damit
\[
  R^2 = 1 - \frac{11}{12{,}67} \approx 1 - 0{,}87 = 0{,}13.
\]
Das Modell erklärt hier etwa 13\,\% der Gesamtvarianz.

\paragraph{Kritische Einordnung}
\(R^2\) ist anschaulich und modellunabhängig, kann jedoch durch Ausreißer oder
sehr homogene Daten verzerrt sein und steigt mit zunehmender Modellkomplexität
(vgl. \cite{Miller2024}).

\chapter{Fazit}
In der Arbeit sollten die Evaluationsmetriken für die ML-Algorithmen zugänglich gemacht werden.
Die Auswahl geeigneter Evaluationsmetriken hängt dabei stark von der Lernaufgabe und
den Zielen ab: Intrinsische Clustering-Indizes (Silhouette, Davies-Bouldin,
Calinski-Harabasz) messen Kohäsion und Separation, während extrinsische Indizes
(Purity, ARI, NMI) die Clusterlösung mit bekannten Klassen vergleichen. In der
Klassifikation sind je nach Datenlage verschiedene Kombinationen von Accuracy,
Precision, Recall, F1 und AUC sinnvoll; keine einzelne Metrik genügt im
Allgemeinen. Für die Regression geben MSE/RMSE, MAE und $R^2$ unterschiedliche
Blickwinkel auf die Fehler. Ein umfassendes Bild entsteht erst, wenn man
mehrere komplementäre Metriken betrachtet.

Wichtig ist, dass man die Stärken
und Schwächen jeder Metrik kennt: Manche Metriken belohnen viele kleine Cluster
(Purity), andere berücksichtigen Zufallseffekte (ARI), und wieder andere
gewichten Ausreißer unterschiedlich (MSE vs. MAE). Letztlich müssen Metriken im Kontext der Anwendung interpretiert werden.

Die einzelnen Modelle sollten unter verschiedenen Blickwinkeln
bewertet und die Ergebnisse kritisch hinterfragt werden. Durch die Kombination
mehrerer Metriken und bestmögliche Verwendung des verfügbaren Expertenwissens
lässt sich die Qualität von Clustering, Klassifikation und Regression im
maschinellen Lernen am aussagekräftigsten beurteilen. 





% ---- Literaturverzeichnis
\cleardoublepage
\renewcommand*{\chapterpagestyle}{plain}
\pagestyle{plain}
\pagenumbering{Roman}                   % Römische Seitenzahlen
\setcounter{page}{\numexpr\value{savepage}+1}
\printbibliography[title=Literaturverzeichnis]

% ---- Anhang
\appendix
%\clearpage
%\pagenumbering{Roman}  % römische Seitenzahlen für Anhang
\chapter{Anhang}

\newpage
\end{document}
